\section{The complete original GCT-IV module and its integral projection}
\label{sec:full-source-generators}

Let $A=\mathbb Z[q,q^{-1}]$, $K=\mathbb Q(q)$, $p=q+q^{-1}$,
and $[j]=q^{j-1}+q^{j-3}+\cdots+q^{1-j}$ for $j>0$.
This section constructs the four original Chevalley matrices on
$\check X_{(7,5,3)}$, rather than extending the four-dimensional
negative-Borel quotient of Section~\ref{sec:gct-generator-interval}.
The column shape is exactly $\alpha=(3,3,3,2,2,1,1)$; the row
shape is $(7,5,3,0)$; the column-count shape is $[0,3,2,2]$.
The source parameter, denoted $u$ in its TeX, remains the
indeterminate $q$ here.

The construction is finite and explicit.  The file
\texttt{agents/}\allowbreak
\texttt{full\_source\_generators/}\allowbreak
\texttt{full\_generators.json}
contains every entry of all four $1260\times1260$ matrices and
the image of every original tabloid. Its coefficient convention is
a list of pairs $(k,a_k)$ denoting the Laurent polynomial
$\sum_k a_kq^k$; its sparse matrix convention lists target indices
in each source column. Indices start at zero. The file
\texttt{source\_basis/basis\_1260.json} gives every original
column label, scalar, invariant record, twelve-tuple, and weight.
The definitions and recurrences below specify these files exactly;
the attached Python programs implement those definitions with
integer arithmetic and exhaustive, finite checks.

\subsection{The original ambient basis and all adjacent relations}

Write $Y=\check Y_\alpha$. Its original NST basis is indexed by
the set $\mathcal W$ of all column tuples in
\[
 \{123,124,134,234\}^3\times
 \{12,13,23,32,24,34\}^2\times\{1,2,3,4\}^2.
\]
The product is the source based product, with the indicated column
composition retained. Thus $|\mathcal W|=4^3 6^2 4^2=36864$.
These are column labels, not ordinary tensor monomials. In
particular, the two columns $23$ and $32$ are distinct.
The alphabet is $1=(1,1),2=(1,2),3=(2,1),4=(2,2)$.
For diagram calculations read each column from bottom to top,
then extract its $V$ and $W$ letters:
\[
\begin{array}{c|cc@{\qquad}c|cc}
c&V&W&c&V&W\\\hline
1&1&1&2&1&2\\3&2&1&4&2&2\\
12&11&21&13&21&11\\23&21&12&32&12&21\\
24&21&22&34&22&21\\
123&211&121&124&211&221\\
134&221&211&234&221&212
\end{array}
\]
Some selected canonical tensor summands of the height-three
columns use other words. The displayed dictionary is the literal
diagram reading rule; it avoids identifying a whole column with
one summand of its canonical expansion.

Let $R_K$ be the original sum of the six adjacent source
submodules $\check Y_{\triangleright^i\alpha}\subset Y$.
The following tables give all of its integral two-column basis
relations. Each row is an equality of the images in $Y/R_K$;
the equality means that the difference belongs to the indicated
local submodule before taking the quotient:
\[
\begin{array}{c|l}
(2,1)&(23)(1)=(32)(1),\quad(23)(3)=(34)(1),\\
 &(24)(3)=(34)(2),\quad(24)(1)=(32)(2)\\\hline
(3,2)&(124)(13)=-(134)(12),\quad(234)(23)=-(234)(32),\\
 &(234)(13)=-(134)(32),\quad(124)(23)=-(234)(12)\\\hline
(2,2)&(32)(24)=(24)(32),\quad(12)(13)=(13)(12),\\
 &(34)(24)=(24)(34),\quad(12)(23)=(23)(12),\\
 &(32)(13)=(13)(32),\quad(34)(13)=(13)(34),\\
 &(34)(23)=(23)(34),\quad(12)(24)=(24)(12),\\
 &(32)(23)=(23)(32),\quad(24)(13)=(34)(12).
\end{array}
\]
The first nine $(2,2)$ relations are its nine-dimensional
degree-preserving component; the last is its one-dimensional
invariant component. For an integral local relation the source
based products for the submodule and ambient module agree, so
these relations hold with every left and right NST context.

There are two further local strings at each equal height. The
active side, its three words $C_{11},C_{12},C_{22}$, and its
invariant $I$ are as follows:
\[
\begin{array}{c|c|ccc|c}
h&\text{side}&C_{11}&C_{12}&C_{22}&I\\\hline
1&V&(2)(1)&(2)(3)&(4)(3)&(4)(1)\\
1&W&(3)(1)&(3)(2)&(4)(2)&(4)(1)\\
2&V&(32)(12)&(32)(32)&(34)(32)&(34)(12)\\
2&W&(23)(13)&(23)(23)&(24)(23)&(24)(13)\\
3&V&(124)(123)&(124)(134)&(234)(134)&(234)(123)\\
3&W&(134)(123)&(134)(124)&(234)(124)&(234)(123)
\end{array}
\]
In the ambient two-column module their projected string vectors
are $\widetilde C_{11}=C_{11}$,
$\widetilde C_{12}=C_{12}+p^{-1}I$,
$\widetilde C_{22}=C_{22}$.  Their contextual products in
$\check Y_{\triangleright^i\alpha}$ give exactly
\begin{equation}\label{eq:full-contextual-relations}
 BC_{11}D=-p^{-1}B'ID,\qquad
 BC_{12}D=-p^{-1}BID,\qquad
 BC_{22}D=-p^{-1}BID'.
\end{equation}
Here is the full meaning of $B'$ and $D'$. Pair every $2$ with
the next available $1$ to its right, using a stack of unmatched
$2$ positions. For $11$, take the leftmost local unpaired $1$.
If it has no partner in $B$, set $B'=0$; otherwise change its
partner from $2$ to $1$ and apply the raising column Kashiwara
operator at that column. For $22$, take the rightmost local
unpaired $2$. If it has no partner in $D$, set $D'=0$;
otherwise change that partner from $1$ to $2$ and apply the
lowering column Kashiwara operator there. In either case an
extra internal arc makes the corrected word zero. The other
context is unchanged. The middle equality changes no context.
These are the source contextual projection formulas, with every
outside letter retained. An isolated equality from this table
cannot be substituted in a longer product while discarding
$B'$ or $D'$.

The integral table and these six strings exhaust the local
submodule bases for the adjacent shapes occurring in $\alpha$:
$(3,3)$ twice, $(3,2)$ once, $(2,2)$ once, $(2,1)$ once and
$(1,1)$ once. Taking all contexts therefore spans exactly $R_K$,
by its defining sum. This proves that the following relation
construction presents the original $\check X_{(7,5,3)}$.

\subsection{A finite exact quotient algorithm}

Order $\mathcal W$ lexicographically using the column orders in
its displayed Cartesian product. There are $22528$ contextual
instances of integral table rows and $47616$ instances of
\eqref{eq:full-contextual-relations}. Each says
$t_i=\epsilon p^kt_j$ with $\epsilon\in\{1,-1\}$ and
$k\in\{0,-1\}$, or $t_i=0$.
Join the vertices by these equalities, keeping for each vertex
an accumulated scalar $\epsilon p^k$ relative to a chosen root.
An equality within one component either agrees with its
accumulated scalar, or yields
$(1-\epsilon p^k)t_{\rm root}=0$. In the latter case set that
component to zero; do the same when a zero relation is present.

This is Gaussian elimination specialized to binomial relations
over $K$. Indeed a spanning tree in each component expresses
every vertex as its root times a nonzero scalar. Every remaining
edge is then either $0=0$ or a nonzero scalar times that root.
Consequently a component contributes exactly one free $K$
coordinate if every cycle agrees and no zero relation is present,
and contributes zero otherwise. For this original presentation
the algorithm leaves exactly $1260$ nonzero components; $18652$
original words are zero and $18212$ original words have nonzero
image. The program checks every equation again through its
computed projection.

This elimination takes place over $K$. Inverting a nonzero
$1-\epsilon p^k$ in a conflicting component does not prove that
an integral presentation obtained merely by clearing denominators
has no torsion. We make no such claim. The integral lattice and
its exact kernel are constructed separately below inside the
original field quotient.

\subsection{The complete original honest basis and its signs}
\label{fullbasis:enumeration}

We specify the original basis at the original parameter.  Put
$K=\mathbb Q(q)$, $p=q+q^{-1}$, $\nu=(7,5,3,0)$, and
$\alpha=\nu'=(3,3,3,2,2,1,1)$.  In this subsection a string such as
$124$ denotes a single column, read from top to bottom; parentheses list
columns from left to right.  In particular $23$ and $32$ are distinct
nonstandard columns.  The source column alphabets, in their source order,
are
\[
\mathcal C_3=(123,124,134,234),\qquad
\mathcal C_2=(12,13,23,32,24,34),\qquad
\mathcal C_1=(1,2,3,4).
\]
All references to source lines below refer to the retained original
\texttt{Apr13KroneckerGCT4.tex}, whose SHA-256 is
\texttt{5c5b990306c3b817d822df9f5088fdd42}\allowbreak
\texttt{dda146fc3449155a0b2f6202d2ed49e}.
The alphabets and their full canonical tensor expressions occur at
lines~4630--4635; their order occurs at~6064.

The ambient source space is
\[
\check Y_\alpha=
 \bigotimes_{j=1}^{7}\check\Lambda^{\alpha_j}X,
\qquad \dim_K\check Y_\alpha=4^3\,6^2\,4^2=36864.
\]
Its NST basis uses the source projected canonical product: if
$T=(C_1,\ldots,C_7)$, its vector is
$\check p_\alpha(C_1\mathbin\heartsuit\cdots\mathbin\heartsuit C_7)$.
Here $\heartsuit$ concatenates the original upper canonical words in each
of the $V,W$ factors, and $\check p_\alpha$ is the source projector.
This is the definition at~4751--4765 and agrees with the based tensor
product by~4795--4805.  Thus these words do not mean an unqualified pure
tensor in the NSC basis.  The original quotient is
\[
 \check X_\nu=\check Y_\alpha/\sum_{j=1}^{6}\check Y_{>^j\alpha},
 \qquad
 \check Y_{>^j\alpha}=
 \check Y_{(\alpha_1,\ldots,\alpha_{j-1})}\otimes
 \check Y_{>(\alpha_j,\alpha_{j+1})}\otimes
 \check Y_{(\alpha_{j+2},\ldots,\alpha_7)}.
\]
These are exactly the objects in source~4918--4920, with the source local
relation submodules.  When written in the canonical NST product the
nonintegral relations must retain their context, as in~5773--5785.

Here is an exhaustive finite index set for the honest basis.  Choose
$i_3,i_2,i_1\in\{0,1\}$ and nonnegative integers $m_C$ for each column
$C\in\mathcal C_3\sqcup\mathcal C_2\sqcup\mathcal C_1$, subject to
\begin{equation}\label{fullbasis:counts}
 \sum_{C\in\mathcal C_3}m_C=3-2i_3,\qquad
 \sum_{C\in\mathcal C_2}m_C=2-2i_2,\qquad
 \sum_{C\in\mathcal C_1}m_C=2-2i_1.
\end{equation}
Retain precisely the tuples satisfying all the following source
conditions:
\begin{align}
 m_2m_3&=0,&
 m_1>0&\ \Longrightarrow\ m_{34}=m_{24}=m_{32}=0,\nonumber\\
 m_2>0&\ \Longrightarrow\ m_{34}=0,&
 m_3>1&\ \Longrightarrow\ m_{24}=0,\nonumber\\
 m_{23}&\leq1,&m_{32}&\leq1,\nonumber\\
 m_{124}m_{134}&=0,&
 m_{234}>0&\ \Longrightarrow\ m_{12}=m_{13}=m_{32}=0,\nonumber\\
 m_{134}>0&\ \Longrightarrow\ m_{12}=0,&
 m_{124}>1&\ \Longrightarrow\ m_{13}=0,\nonumber\\
 m_{234}m_1&=0.&&\label{fullbasis:predicates}
\end{align}
Within each height the invariant-free columns are in the displayed
alphabet order.  Equations~\eqref{fullbasis:predicates} retain every
condition of source~6114--6130; conditions with identical antecedents
have been grouped without changing their content.

From the tuple form the exact source word
\begin{equation}\label{fullbasis:word}
\begin{aligned}
 T(m,i)={}&123^{m_{123}}124^{m_{124}}134^{m_{134}}
 (234,123)^{i_3}234^{m_{234}}\\[-1mm]
 &12^{m_{12}}13^{m_{13}}23^{m_{23}}32^{m_{32}}
 (24,13)^{i_2}24^{m_{24}}34^{m_{34}}\\[-1mm]
 &1^{m_1}2^{m_2}3^{m_3}(4,1)^{i_1}4^{m_4}.
\end{aligned}
\end{equation}
Exponentiation of a column or pair means that many consecutive copies.
The invariant insertion positions and orientations are those of
source~6169.  The invariant record is $(0,i_3,i_2,i_1)$ and the degree
is $d=i_3+i_2+i_1$.  Indeed the inserted pairs have two arcs connecting
their columns and no arcs exiting the pair.  Removing them gives the
invariant-free word with multiplicities $m_C$.  The source classification
at~6263--6269 identifies these straightened words bijectively with the
honest classes.  The complete quotient computation in this chapter also
checks, independently for this fixed shape, that the words in
\eqref{fullbasis:word} project to distinct nonzero components and exhaust
all nonzero quotient components.  Thus both the representatives and the
module are the original source objects.

For completeness we give an exact diagram algorithm.  Read the columns
bottom to top and then left to right and apply
\[
 1\mapsto(1,1),\quad2\mapsto(1,2),\quad
 3\mapsto(2,1),\quad4\mapsto(2,2).
\]
This produces the original length-$15$ $V$ and $W$ words (source~4723).
For each word scan left to right, pushing every $2$ on a stack; a $1$
pairs with the last unmatched $2$ if there is one.  Otherwise it remains
unpaired.  This is exactly the source parenthesis rule~2448.  Record the
column containing each letter.  An arc is external when its ends lie in
different columns; two arcs between a pair of columns give an invariant
pair (source~4938--4954).  A two-edge component involving one column of
each height $3,2,1$ is a $3$--$2$--$1$ arc.  Its color is the common
edge color, or the color of the longer edge if the colors differ, as in
source~6320.  All other remaining edges are pure.  A height-$2$ column
with no external edge is externally free; its type is $V$ for
$12,32,34$ and $W$ for $13,23,24$ (source~6068--6073).

The exact twelve-entry label is
\begin{equation}\label{fullbasis:zeta}
 \zeta=(i_4,i_3,i_2,i_1,a_{21},a_{32},a_{31V},a_{31W},
                 a_{321V},a_{321W},f_V,f_W).
\end{equation}
Here $a_{21},a_{32}$ count pure arcs, the other $a$'s retain their
specified source colors, and $f_V,f_W$ count externally free columns.
This is the original order in source~7233--7238.  Together with the two
unpaired words, necessarily of the form $1^a2^b$, it determines the
honest class by~6473--6481.  Original torus weights count all letters,
including paired letters: each factor's weight is
$(\#1,\#2)$.  This retains the determinant-weight contribution of every
invariant.

\paragraph{A source sign ambiguity and its exact resolution.}
Source~5489--5490 specifies a positive class using the multiplier
$(-1/p)^d$ and a representative with a $W$-colored $3$--$2$ arc whenever
the class has a $3$--$2$ arc.  It also asserts that having no $3$--$2$
arc for every representative is equivalent to having none for any
representative.  That latter assertion fails; the $3$--$2$--$1$ example
at~6337--6388 already has equivalent representatives of both types.
We preserve the original sign choice through its stated
$W$-arc representative criterion and explicitly record the needed
correction, instead of treating the inaccurate equivalence as a theorem.

Here is the exact integral witness relevant to this shape:
\begin{equation}\label{fullbasis:signwitness}
 \begin{aligned}
 (134,13,23,1,x)
 &=(134,13,32,1,x)\\
 &=(134,32,13,1,x)\\
 &=-(234,13,13,1,x),\qquad x\in\{1,2,3,4\}.
 \end{aligned}
\end{equation}
The first equality is the source integral $(23,1)=(32,1)$ relation
at~5080.  The second is the type-$V$/type-$W$ commutation at~5139.
The last is $(134,32)=-(234,13)$ at~5209.  All are degree-preserving
integral moves, so their insertion into the retained contexts is valid.
The original left word has a $3$--$2$--$1$ $V$-arc and no $3$--$2$
edge; the right word has a $W$-colored $3$--$2$ edge.  Direct application
of the pairing algorithm verifies both assertions for each of the four
values of $x$.  Consequently the positive vector represented by the
left word is $-(-1/p)^d$ times that word.

Within the present shape, the additional words requiring this sign are
exactly
\[
\begin{gathered}
 (123,123,134,13,23,1,x),\quad
 (123,134,134,13,23,1,x),\\
 (134,134,134,13,23,1,x),\quad
 (134,234,123,13,23,1,x),\\
 x\in\{1,2,3,4\}.
\end{gathered}
\]
For the last family first use
$(134,234,123)=(234,123,134)$, the invariant commutation
at~5286--5287, then apply~\eqref{fullbasis:signwitness}.  The first
three families use~\eqref{fullbasis:signwitness} directly in their
height-$3$ prefixes.  Exhaustiveness follows by evaluating
\eqref{fullbasis:predicates} and the pairing algorithm on the finite
index set; every index and its predicate values are computed by
\texttt{enumerate\_basis.py}, and independently by
\texttt{audit/independent\_audit.py}.  All sixteen equalities are also
checked against the full original quotient projection by
\texttt{check\_sign\_witnesses.py}.  In source lexicographic order their
zero-based indices are $262$--$265$, $478$--$481$, $1000$--$1003$, and
$1108$--$1111$.

Thus the final completely explicit scaling is
\begin{equation}\label{fullbasis:sign}
 b_{m,i}=\epsilon(T)(-1/p)^d T,
 \qquad
 \epsilon(T)=
 \begin{cases}
 -1,&\text{$T$ has a $V$-colored $3$--$2$ edge},\\
 -1,&\text{$T$ has no $3$--$2$ edge and has a $3$--$2$--$1$ $V$-arc},\\
 1,&\text{otherwise}.
 \end{cases}
\end{equation}
The sixteen corrections affect neither the index set nor any original
quotient relation.  They give an explicitly specified diagonal sign
change, with its exact inverse equal to itself.  All four retained
labels below have their original established scalings unchanged.

\paragraph{Exhaustive enumeration and a second exact bijection check.}
For each invariant record the numbers of words in
\eqref{fullbasis:word} are as follows:
\[
\begin{array}{c|rrrrrrrr}
(i_3,i_2,i_1)&000&001&010&011&100&101&110&111\\\hline
\text{number}&628&148&108&16&268&56&32&4.
\end{array}
\]
Their sum is $1260$; the degree counts are $628,524,104,4$ for
$d=0,1,2,3$.  All $1260$ keys consisting of~\eqref{fullbasis:zeta}
and the two unpaired words are distinct.  There are $48$ distinct
twelve-tuples.  For each such tuple, if its two unpaired lengths are
$h_V,h_W$, its vectors contain exactly all pairs
\[
 (1^{h_V-a}2^a,1^{h_W-b}2^b),\qquad
 0\leq a\leq h_V,\quad0\leq b\leq h_W.
\]
The complete data are in \texttt{basis\_1260.json} and
\texttt{ssyt\_bijection\_check.json}.

We additionally verify the original map to ordinary semistandard
tableaux, retaining every source case rather than using its cardinality
alone.  For clarity here is that map in multiplicity form.  Let $m'_C$
be the ordinary output multiplicities, initially equal to the
invariant-free $m_C$; set $m'_{32}=0$.  The columns $123,234,1,4$ keep
their invariant-free multiplicities.  The following tables specify
all other multiplicities; entries not changed in a row keep their
invariant-free values.  First, the right part (source~6195--6210):
\[
\begin{array}{c|cc|cc}
\text{condition}&m'_2&m'_3&m'_{24}&m'_{34}\\\hline
m_1>0&i_1+m_2&i_1+m_3&m_{24}&m_{34}\\
m_1=0,\ m_{24}>0,\ m_2=0&0&2i_1+m_3&m_{24}&m_{34}\\
m_1=0,\ m_{24}>0,\ m_2>0&i_1+m_2&i_1&m_{24}&0\\
m_1=m_{24}=0,\ i_1>0&i_1+m_2&i_1+m_3&m_{34}&0\\
m_1=m_{24}=i_1=0&m_2&m_3&0&m_{34}
\end{array}
\]
Next the middle part (source~6217--6232):
\[
\begin{array}{c|cc}
\text{condition}&m'_{14}&m'_{23}\\\hline
m_{234}=0,\ m_1>0&2i_2+m_{23}&0\\
m_{234}>0,\ m_1=0&0&2i_2+m_{23}\\
m_{234}=m_1=m_{23}=0&2i_2+m_{32}&0\\
m_{234}=m_1=0,\ m_{23}>0&0&2i_2+m_{32}+m_{23}
\end{array}
\]
Finally the left part (source~6237--6250):
\[
\begin{array}{c|cc|cc}
\text{condition}&m'_{124}&m'_{134}&m'_{12}&m'_{13}\\\hline
m_{234}>0&i_3+m_{124}&i_3+m_{134}&m_{12}&m_{13}\\
m_{234}=0,\ m_{13}>0,\ m_{134}=0&2i_3+m_{124}&0&m_{12}&m_{13}\\
m_{234}=0,\ m_{13}>0,\ m_{134}>0&i_3&i_3+m_{134}&0&m_{13}\\
m_{234}=m_{13}=0,\ i_3>0&i_3+m_{124}&i_3+m_{134}&0&m_{12}\\
m_{234}=m_{13}=i_3=0&m_{124}&m_{134}&m_{12}&0
\end{array}
\]
Conditions~\eqref{fullbasis:predicates} imply that exactly one row of
each table applies.  The three tables write disjoint groups of output
multiplicities and use only original $m_C$, so they do not interfere.
Arrange the output columns in ordinary column lexicographic order.

The exact finite verification in \texttt{check\_ssyt\_bijection.py}
enumerates every tuple in~\eqref{fullbasis:counts} by combinations with
repetition, tests every predicate~\eqref{fullbasis:predicates}, and
evaluates the three tables.  Independently it enumerates every ordinary
height-$3$ column multiset of size $3$, height-$2$ column multiset of
size $2$, and height-$1$ column multiset of size $2$.  Each ordinary
column is a strictly increasing sequence in $\{1,2,3,4\}$; every common
row coordinate of adjacent columns is tested for weak increase.  These
are exactly the defining semistandard conditions, and the enumeration
has no other exclusion.  The two sets are equal, have $1260$ elements,
and the map's $1260$ outputs are distinct.  The complete inverse is
therefore given without an existence assumption by the reversed
finite correspondence in \texttt{ssyt\_correspondence.json}.
As a further exact check the ordinary Weyl product is
\[
 \prod_{1\leq i<j\leq4}\frac{\nu_i-\nu_j+j-i}{j-i}
 =3\cdot3\cdot\frac{10}{3}\cdot3\cdot\frac72\cdot4=1260.
\]
The numerical identity is a check on the explicitly constructed
bijection, not a replacement for its construction.

\paragraph{Retained original labels.}
In the zero-based basis order defined above,
\begin{align*}
 b_{126}=T&=(123,123,124,12,12,1,1),\\
 b_{146}=T^1&=(123,123,124,23,12,1,1),\\
 b_8=T^2&=-p^{-1}(123,123,123,12,12,4,1),\\
 b_{144}=Q&=-p^{-1}(123,123,124,13,12,4,1).
\end{align*}
For $T^1$ and $Q$, the lexicographic word has the $12$ column before
$23$ and $13$, respectively; their equality with the displayed original
representatives is the integral type-$V$/type-$W$ commutation.  Their
twelve-tuples and unpaired lengths are
\[
\begin{array}{c|c|c}
T&(0,0,0,0,0,0,0,1,0,0,2,0)&(9,3)\\
T^1&(0,0,0,0,0,0,0,0,0,1,1,0)&(7,3)\\
T^2&(0,0,0,1,0,0,0,0,0,0,2,0)&(7,3)\\
Q&(0,0,0,1,0,1,0,0,0,0,1,0)&(5,3).
\end{array}
\]
All these unpaired words consist entirely of $1$'s.  The original
weights $(\mathrm{wt}_V,\mathrm{wt}_W)$ are respectively
$((12,3),(9,6))$, $((11,4),(9,6))$, $((11,4),(9,6))$,
and $((10,5),(9,6))$.

The literal source target at~7333 is retained separately:
\[
 b_{154}=T'_{\mathrm{printed}}
 =-p^{-1}(123,123,124,23,12,4,1).
\]
It has the same twelve-tuple as $Q$ but unpaired words $(11111,112)$
and original weights $((10,5),(8,7))$.  Since $F_V$ commutes with
the $W$ torus, its coefficient in $F_V^2T$ or $F_V^{(2)}T$ is zero.
This is the precise weight obstruction to the printed target; the
corrected $Q$ retains the already proved contextual endpoint and is
not silently identified with the printed object.

\paragraph{Original ambient generator algorithm.}
For $Z=V,W$, let $\varphi_Z(T)$ count its unpaired $1$'s.  Source
4881--4908 gives
\[
 F_ZT=\sum_{j=1}^{\varphi_Z(T)}[j]\,F_{(j)Z}T.
\]
The $j$th operation changes the column containing the $j$th unpaired
$1$, in left-to-right order, by its local Kashiwara arrow; it is zero
if changing that selected letter creates an extra internal arc.  For
$E_Z$, replace the $j$th unpaired $2$ counted from the right, with
coefficient $[j]$, exactly as in source~2508--2517.  The local $F$
strings are
\[
\begin{array}{c|ccccc}
V&1\to3&2\to4&123\to134&124\to234&12\to32\to34\\
W&1\to2&3\to4&123\to124&134\to234&13\to23\to24.
\end{array}
\]
All unlisted arrows are zero, and the local $E$ arrows reverse these
strings.  Choosing the first of the two unpaired $1$'s in a height-$2$
spin-one column creates an extra internal arc; the second gives its
first $F$ arrow.  This zero test is performed for each summand before
the contextual quotient reduction.  Scaling the input by
\eqref{fullbasis:sign}, reducing in the original quotient, and expressing
each output in that same signed basis specifies every original
Chevalley coefficient.  The complete matrices, divided powers, and
their exact checks are supplied by the full quotient construction in
this chapter.


For each of the $1260$ displayed representatives $b_n$, let
$b_n=\sigma_n p^{-d_n}\overline t_{w_n}$, with its fully
specified sign and degree. Each nonzero relation component
contains exactly one such representative. This is checked by
matching all $1260$ roots; hence the vectors $b_n$ are an actual
basis of the original quotient. For every original word $w$
the table records
\begin{equation}\label{eq:full-tab-projection}
 \pi(t_w)=0\quad\hbox{or}\quad
 \pi(t_w)=\epsilon_w p^{k_w}b_{n_w}.
\end{equation}
The algorithm produces only $k_w\in\{0,1,2,3\}$.
The numbers of nonzero words with these four exponents are
$12544,5100,552,16$, respectively. No sign is suppressed in
the table.

\subsection{Every entry of all four original generators}

For a side $Z\in\{V,W\}$ concatenate the displayed diagram
words of an original tabloid $w$. Let its unpaired ones from
left to right be $a_1,\ldots,a_\varphi$, and its unpaired twos
from left to right be $c_1,\ldots,c_\varepsilon$. Define
$\mathcal F_{jZ}(w)$ by changing the column containing $a_j$
with its lowering column Kashiwara operator; define
$\mathcal E_{jZ}(w)$ by raising the column containing
$c_{\varepsilon-j+1}$. An extra internal arc gives zero.
To specify the column operator without an implicit convention:
the other side's weight and internal-arc number stay fixed,
the chosen side's number of twos changes by one, and the chosen
side's original internal-arc number stays fixed. Among the
columns of that same height there is exactly one such column.
If the intermediate changed diagram has more internal arcs
than the old diagram, the result is zero instead.

The original ambient action is exactly
\begin{equation}\label{eq:full-original-action}
 F_Zt_w=\sum_{j=1}^{\varphi_Z(w)}[j]t_{\mathcal F_{jZ}(w)},
 \qquad
 E_Zt_w=\sum_{j=1}^{\varepsilon_Z(w)}[j]t_{\mathcal E_{jZ}(w)}.
\end{equation}
Terms whose changed word is zero are omitted. Apply
\eqref{eq:full-tab-projection} to every summand and multiply by
$\sigma_np^{-d_n}$ to obtain the column for $g b_n$.
Equivalently, for $g=F_Z$ or $E_Z$ its $(m,n)$ entry is the
finite sum
\begin{equation}\label{eq:full-all-matrix-entries}
 (G_g)_{mn}=
 \sum_{\substack{j:\,\pi(t_{g_{jZ}(w_n)})
       =\epsilon_jp^{k_j}b_m}}
       [j]\sigma_n\epsilon_jp^{k_j-d_n}.
\end{equation}
This supplies every matrix entry, including all cancellations,
signs, and possible intermediate denominators.

The nonzero entry counts are
\[
\begin{array}{c|rrrr}
g&F_V&E_V&F_W&E_W\\\hline
\#\operatorname{supp}(G_g)&1998&2083&1998&2083.
\end{array}
\]
Every resulting entry
lies in $A$: finite integral Laurent addition and division by
$p$ with remainder zero verify this directly. The original
projection commutes with the action on every one of the $36864$
ambient words for each of the four generators, including all
$18652$ words with zero image. These $147456$ equalities prove
descent throughout the original relation space, independently
of the choice of representatives.

Let $v_n$ and $w_n'$ be the numbers of $V$ and $W$ ones in
the fifteen-letter diagram of $b_n$. The four diagonal
$\mathfrak{gl}_2$ weight generators have eigenvalues
$q^{v_n},q^{15-v_n},q^{w_n'},q^{15-w_n'}$. Thus the two
Cartan ratios are
$K_Vb_n=q^{2v_n-15}b_n$ and
$K_Wb_n=q^{2w_n'-15}b_n$; their central products are $q^{15}$.
The exact matrices satisfy
\[
 [E_Z,F_Z]=\frac{K_Z-K_Z^{-1}}{q-q^{-1}},\qquad
 [g_V,h_W]=0\quad(g,h\in\{E,F\}),
\]
and $K_ZE_ZK_Z^{-1}=q^2E_Z$,
$K_ZF_ZK_Z^{-1}=q^{-2}F_Z$, with trivial cross Cartan
commutators. These equalities are checked as equality of
integer Laurent coefficient dictionaries in every matrix
column. They also follow by descent from
\eqref{eq:full-original-action}, the source graphical action
on the original tensor product. Thus these are the full
original $U_q(\mathfrak{gl}_2\oplus\mathfrak{gl}_2)$
generators, with their exact quotient identification.

\subsection{Integral lattices, divided powers, and specialization}

Define $L_A=\bigoplus_{n=0}^{1259}Ab_n\subset Y/R_K$.
The exact Laurent integrality just proved makes this an
integral generator lattice. For each $g\in\{E_V,F_V,E_W,F_W\}$
define the matrices recursively by
\[
 G_g^{(0)}=I,\qquad G_g^{(m)}=G_gG_g^{(m-1)}/[m].
\]
This retains the original divided-power denominator until the
division is established. The certificate computes each
coefficient by exact integral Laurent division and verifies
zero remainder for every $m=1,\ldots,10$. Both raising and
lowering matrices have $G_g^{(9)}\ne0$ and $G_g^{(10)}=0$.
Consequently all divided powers preserve $L_A$ and every
higher divided power is zero. The nonzero-entry counts are
\[
\begin{array}{c|rrrrrrrrrrr}
m&0&1&2&3&4&5&6&7&8&9&10\\\hline
F_V,F_W&1260&1998&2314&2100&1489&892&447&185&56&12&0\\
E_V,E_W&1260&2083&2425&2082&1486&907&452&186&56&12&0.
\end{array}
\]
Over $K$ the recurrence is $G_g^m/[m]!$, by induction.
The divided-power identities therefore hold over $K$; both
sides are integral matrices, and the injection $A\hookrightarrow K$
proves the same identities over $A$. This proves the full
divided-power integral action, not only the four ordinary
generator identities.

There is also an exact map from the original ambient integral
word span. Put $Y_A=\bigoplus_{w\in\mathcal W}At_w$ and
use \eqref{eq:full-tab-projection} to define
$\pi_A:Y_A\to L_A$. Its exponents are nonnegative, so this
map is defined over $A$. For each $n$ choose the first word
$w(n)$ in the fixed order whose projection is
$\epsilon(n)b_n$, with exponent zero. The exhaustive table
contains such a word for every $n$. Define
$s_A(b_n)=\epsilon(n)t_{w(n)}$. Then
$\pi_As_A(b_n)=b_n$ for each basis vector, proving
$\pi_As_A=I$.

Here is the entire kernel, with no unspecified torsion. For a
nonpivot word $w$, set
\[
 h_w=\begin{cases}
 t_w,&\pi_A(t_w)=0,\\
 t_w-\epsilon_w\epsilon(n_w)p^{k_w}t_{w(n_w)},
       &\pi_A(t_w)=\epsilon_wp^{k_w}b_{n_w}.
 \end{cases}
\]
The $35604$ vectors $h_w$ form an $A$-basis of $\ker\pi_A$.
Their image is zero by the displayed formula. They are
independent because in a relation among them the coefficient
of each distinct nonpivot original word is its own coefficient.
Subtracting $\sum_w a_wh_w$ from any kernel vector leaves
only pivot words; their images are signed distinct basis
vectors of $L_A$, so all remaining coefficients vanish.
This proves spanning. Therefore
\begin{equation}\label{eq:full-integral-split}
 0\longrightarrow A^{35604}\xrightarrow{h}Y_A
 \xrightarrow{\pi_A}L_A\longrightarrow0
\end{equation}
is split exact as a sequence of $A$-modules, with cokernel
zero. Its section is a specified linear section; no
equivariance of that section is asserted. The projection
itself is equivariant by the $147456$ descent identities.
The ambient graphical module has its original divided-power
action, so its kernel is stable under that action as well.
The section and all kernel vectors are listed in
\texttt{raw\_integral\_projection.json}.

Upon $q=1$ over $\mathbb Z$, $p=2$, and the same formulas
with $p^{k_w}=2^{k_w}$ remain a split exact sequence of free
abelian groups. In particular the natural kernel base-change
map is an isomorphism. Every full generator and divided-power
matrix specializes by replacing $\sum a_kq^k$ by $\sum a_k$;
\texttt{specialized\_generators.json} gives all these exact
integer entries. Thus this specialization retains every
coefficient and has rank $1260$ over $\mathbb Z$.
This is a statement about the actual image lattice inside
the field quotient. It does not replace the unsaturated
integral span of the original binomial relations by its
saturation without identifying the map.

For clarity the corresponding saturation map can also be
stated exactly. Let $P_A\subset Y_A$ be generated by the
integral contextual rows and by every nonintegral row
$pt_w+t_{w'}$ (or $t_w$ when its correction is zero).
Then $P_A\subset\ker\pi_A$ and
$K\otimes_AP_A=K\otimes_A\ker\pi_A=R_K$.
The induced map $Y_A/P_A\to L_A$ is onto and has exactly
the torsion kernel $(\ker\pi_A)/P_A$.
Here ``torsion'' follows because tensoring that kernel with
$K$ is zero: each of its elements is annihilated by some
nonzero element of $A$. It need not be zero. A fully explicit
presentation of this kernel is obtained by expressing each
displayed generator of $P_A$ in the basis $h_w$: its
coefficient at $h_w$ is just its coefficient at that nonpivot
word, as proved in the elimination above. This preserves
the original relation presentation and states the exact
map to its torsion-free image lattice.

There is a sharper exact answer after a specified
localization. Put $B=A[p^{-1}]$, retaining $2$ as a
nonunit. In this ring every tree scalar in the original
relation graph is a unit. Elimination along a spanning
tree is therefore an invertible change of presentation
coordinates over $B$. For each connected component it
leaves one root and relations
$(1-\epsilon p^k)\,\mathrm{root}=0$ for its non-tree
edges, plus $\mathrm{root}=0$ if there is an explicit
zero relation. The exhaustive graph has $1260$ components
with no defect and no zero relation, $4147$ with an
explicit zero relation, and $40$ further components
whose entire nonzero cycle ideal is generated by
$1-(-1)=2$. Consequently the actual localized
presentation has the exact decomposition
\begin{equation}\label{eq:full-localized-two-torsion}
\begin{aligned}
 B^{36864}/P_B&\simeq B^{1260}\oplus(B/(2))^{40},\\
 \ker(B^{36864}/P_B\longrightarrow B\otimes_A L_A)
       &\simeq(B/(2))^{40}.
\end{aligned}
\end{equation}
The map to the image lattice is the previously specified
original-word projection; the displayed isomorphism is
given by the spanning-tree coordinates. This identifies
the kernel of the actual map, with every original
component listed in
\texttt{quotient\_audit/}\allowbreak
\texttt{relation\_graph\_certificate.json}.

One of the forty torsion generators has the explicit
original representative
\[
u=(123)(123)(234)(23)(23)(1)(1).
\]
The three original integral rows give, in succession,
\[
\begin{aligned}
u&=-v,\qquad
 v=(123)(123)(234)(32)(23)(1)(1),\\
v&=w,\qquad
 w=(123)(123)(234)(23)(32)(1)(1),\\
w&=u.
\end{aligned}
\]
The first row uses the adjacent $(3,2)$ relation, the
second commutes $(32)(23)$, and the third uses
$(32)(1)=(23)(1)$. Thus $2u=0$. Its component has
six original vertices, no explicit zero relation, and
cycle ideal exactly $(2)$; the spanning-tree proof shows
that $u\ne0$ and generates one $B/(2)$ summand.
Indeed $B/(2)$ is the localization of
$\mathbb F_2[q,q^{-1}]$ at its nonzero element
$q+q^{-1}$, so it is not the zero ring. This proves
nonvanishing, not merely the absence of a proof that
the torsion class vanishes.

\subsection{The original rational specialization bridge}

We first specify the source's map $s_X$, which is a map of
quotients. For any original NST $t_w$, let $d(w)$ be its
number of invariant column pairs, including pairs of
different heights. Such a pair is detected by two arcs
between its two columns. Define
$Y_A^S=\bigoplus_w A(-p^{-1})^{d(w)}t_w$ inside $Y$.
Let $R_A^S$ be the source's sum of the correspondingly
scaled adjacent submodule lattices: scale each contextual
submodule basis element $\widetilde t$ by
$(-p^{-1})^{\deg(\widetilde t)}$, with its ambient filtered
degree retained. Thus the source has the canonical surjection
\[
 s_X:Y_A^S/R_A^S\longrightarrow\pi(Y_A^S),\qquad
 [y]\longmapsto\pi(y).
\]
The identification of its localized relation lattice
with $P_B$ requires more than equality of their
$K$-spans. We give both containments explicitly.
By the source definition at lines $5620$--$5624$, a
contextual basis vector is
$g=B_0\mathbin{\heartsuit_g}\widetilde C
\mathbin{\heartsuit_g}D_0$, with $B_0,D_0$ ranging
over all original NST bases of the two contexts.
Its actual lattice generator at lines $5881$--$5889$
is $(-p^{-1})^{\deg(g)}g$. For an integral local
$\widetilde C$, the equality at lines $5680$--$5685$
is exact in $Y$: $g$ equals the ambient based product
of the listed binomial with its two contexts. Thus it
is the corresponding integral graph row with its
original coefficients $+1,-1$ or $+1,+1$.
For a nonintegral local $\widetilde C$, the equality
at lines $5773$--$5785$ is also exact in $Y$:
$g=t_w+p^{-1}t_{w'}$, or $g=t_w$ when the contextual
correction is zero. This uses the explicit contextual
corollary, not the weaker congruence for nonintegral
vectors in lines $5680$--$5685$. The cleared graph
row is respectively $pg$ or $g$. Therefore each
source lattice generator is $\pm p^j$ times its
corresponding graph row for a specified integer $j$,
and conversely each graph row is $\pm p^{-j}$ times
that source lattice generator. Every such coefficient
is a unit in $B$; no factor $2$ is inverted in either
containment. The contexts and local rows are
exhaustive in both definitions. It follows that
$B\otimes_A R_A^S=P_B$ as actual submodules of
$B\otimes_A Y_A^S=B^{36864}$.
Hence the localization of this exact source presentation is the presentation
in \eqref{eq:full-localized-two-torsion}; its localized
map $s_X$ has precisely the same nonzero kernel
$(B/(2))^{40}$. In particular the source remark at
lines $5901$--$5908$ asserting equality of its two
integral forms cannot be used as an injectivity proof
for these literal prescribed contextual lattices.
The displayed original cycle gives the exact integral
obstruction, while the rational specialization below
constructs their isomorphism over its stated base.
In this shape $\pi(Y_A^S)=L_A$: applying the explicit
projection to each scaled word gives zero or
$\epsilon(-1)^{d(w)}p^{k_w-d(w)}b_{n_w}$.
All nonzero exponents are $0$ or $1$; their counts are
$17728$ and $484$. Conversely the scaled honest
representatives include every $b_n$. This proves both
containments without assuming that $s_X$ is injective over $A$.

Its rational specialization is injective, by the following
additional argument. Every edge of the original relation
graph has scalar $\epsilon p^k$, and every cycle product
has that form. At $p=2$ all edges remain nonzero, and
$\epsilon 2^k=1$ if and only if $\epsilon=1$ and $k=0$.
These are exactly the conditions for $\epsilon p^k=1$ in
$K$. Hence the $1260$ nonzero components remain nonzero,
and every zero component remains zero. More specifically,
the independently reconstructed graph has $4147$ components
with an explicit zero relation and $40$ further components
killed by a cycle of scalar $-1$. The specialized relation
space has rank $36864-1260=35604$ over $\mathbb Q$.
At $q=1$ the rescalings $(-p^{-1})^d=(-1/2)^d$ are
nonzero rational numbers, so they change no rational span.
Thus $\mathbb Q\otimes_A(Y_A^S/R_A^S)$ has dimension
$1260$. The specialization of $s_X$ is a surjection onto
$\mathbb Q\otimes_A L_A$, also of dimension $1260$ by
\eqref{eq:full-integral-split}; hence it is an isomorphism.
Its formula is the original quotient map, not a dimension
comparison without a specified morphism.

The rational specialization of $L_A$ is the restriction of
the classical Schur module of row shape $(7,5,3,0)$, through
an explicit quotient map, as follows. Realize each original
column at $q=1$ in $(V\otimes W)^{\otimes h}$ by its source
canonical column formula; take the ordinary exterior quotient
$e_h:(V\otimes W)^{\otimes h}\to\bigwedge^h(V\otimes W)$
on each column factor. The restriction to the original
nonstandard exterior subspace is a $\mathrm{GL}(V)\times
\mathrm{GL}(W)$ isomorphism over $\mathbb Q$. To see
injectivity, define $\operatorname{Alt}_h:
\bigwedge^hX\to X^{\otimes h}$ by the unnormalized signed
sum over all permutations. Directly from the alternating
relations, $e_h\operatorname{Alt}_h=h!I$ on $\bigwedge^hX$
and $\operatorname{Alt}_he_h=h!I$ on the alternating
tensor subspace. The source subspace at $q=1$ is exactly
this alternating image, since the two tensor permutation
actions on $V,W$ combine to the permutation action on
$X=V\otimes W$. Thus $e_h$ has inverse
$\operatorname{Alt}_h/h!$ on these subspaces. This proof retains
the factor $h!$ and works over $\mathbb Q$, as stated.

Tensor these seven maps, keeping the original ordered column
factors. Apply this tensor product linearly to the actual
specialized based-product vector $t_w$ and its tensor
expansion. A based product at $q=1$ need not be the ordinary
tensor product of its column labels; replacing each label
independently and concatenating would not specify this map.
The source defines its two-column submodules as
the exact $U^\tau$ summands removed by the classical
two-column Schur quotient: this is its defining prescription
at source lines $5015$--$5021$, not an inference from the
global dimension $1260$. Here the ambient two-column module
is itself multiplicity-free as a $U^\tau$-module; the
involution $\tau$ interchanges $V,W$ and distinguishes
otherwise repeated restrictions. We verify the small cases
to justify uniqueness of the prescribed subspace. Write
$D_V=\det V,D_W=\det W$. First
$Y_{(1,1)}$ has the three distinct $U^\tau$ summands of
dimensions $9,6,1$, obtained from
$(S^2V\oplus D_V)\otimes(S^2W\oplus D_W)$ by pairing
the two cross terms. Next
\[
 \bigwedge^2X=(S^2V\otimes D_W)
                 \oplus(D_V\otimes S^2W).
\]
Tensor with $X$ and use
$S^2V\otimes V=S^3V\oplus D_VV$ and its $W$ version.
The two cross outer terms pair to dimension $16$; the
two copies of $D_VV\otimes D_WW$ split into the two
distinct $\tau$ signs, each of dimension $4$.
Thus $Y_{(2,1)}$ is multiplicity-free with dimensions
$16,4,4$, the two last summands having opposite signs.
For $Y_{(2,2)}$ use
$S^2V\otimes S^2V=S^4V\oplus D_VS^2V\oplus D_V^2$
and its $W$ version in the previous exterior decomposition.
The outer terms pair to dimensions $10$ and $6$;
the two cross terms of dimension $9$ split with opposite
$\tau$ signs, and the two invariant lines also split
with opposite signs. These six summands, of dimensions
$10,6,9,9,1,1$, are pairwise inequivalent as
$U^\tau$-modules. The displayed tensor decompositions
follow by the explicit symmetric product and alternating
contraction maps for two-dimensional $V,W$: in each
case the weights are the disjoint strings with respective
highest weights $(a+b-j,j)$, for $0\le j\le\min(a,b)$;
their dimensions sum to $(a+1)(b+1)$, proving exhaustion.
Finally $\bigwedge^3X\simeq D_VD_W\otimes X$ with
the corresponding overall $\tau$ sign. Hence
$Y_{(3,2)}$ and $Y_{(3,3)}$ are character twists of
$Y_{(1,2)}$ and $Y_{(1,1)}$ and remain
multiplicity-free. This proves the ambient fact needed
here. Selecting the source's specified classical
quotient summands therefore uniquely determines its
kernel as the sum of all remaining ambient summands.
The local bases in the two tables are the explicit bases
given for this prescription in the source's Figures
\texttt{straightening11} through \texttt{straightening33}.
Consequently the columnwise exterior maps take each such
specialized local submodule to the classical adjacent
straightening submodule. This is the exact local
identification used in the source proof at lines
$6013$--$6018$; all its contexts are the unchanged tensor
factors. Quotient by their sum on both sides. The relation
graph argument above identifies the specialized relation
space with the entire specialized kernel, so the resulting
map is
\[
 s_{(7,5,3)}:\mathbb Q\otimes_A L_A
   \longrightarrow
 \left(\bigwedge^3X\right)^{\otimes3}
 \otimes\left(\bigwedge^2X\right)^{\otimes2}
 \otimes X^{\otimes2}\big/\text{classical adjacent relations}.
\]
Its formula takes any representative in $Y_A$ to the class
of its columnwise exterior images. Two representatives
differ by the kernel, whose rational specialization was
proved equal to the specialized relation space; hence the map is
well defined and injective. It is onto because the
columnwise exterior maps are onto. The target is the
classical Schur module $S_{(7,5,3)}X$ by its exterior-column
presentation, and each map is equivariant. This is the
explicit rational specialization correspondence used for
the source character calculation. No literal identification
of a crystal-highest tabloid with an $E$-annihilated vector
is used here.

\subsection{The full source map to the retained interval}

In the full basis the original vectors are exactly
\[
 T=b_{126},\qquad T_1=b_{146},\qquad T_2=b_8,\qquad Q=b_{144}.
\]
These equalities include their original signs and the factors
$-p^{-1}$ in $T_2,Q$. The height-two columns in $T_1,Q$
commute by an integral table row when their order differs
from the representative table; no scalar correction occurs.

Form the directed support graph of the full $F_V,F_W$
matrices. Let $\mathcal D$ be the descendants of $126$,
and let $\mathcal J$ be the descendants which cannot reach
$144$. Exhaustive traversal of these explicit matrices gives
\[
 |\mathcal D|=420,\quad |\mathcal J|=416,\quad
 \mathcal D\setminus\mathcal J=\{126,146,8,144\}.
\]
Both sets are closed under outgoing lowering edges: a
successor remains a descendant, and a successor able to
reach $144$ would provide such a path from its predecessor.
Their $A$-spans are therefore stable under the integral
negative-Borel action and its divided powers. The exact
projection $\bigoplus_{n\in\mathcal D}Ab_n\to M_A$ kills
the $416$ basis vectors in $\mathcal J$ and takes the four
remaining vectors, in the displayed order, to
$e_0,e_1,e_2,e_3$. Its kernel is exactly the displayed
$416$-dimensional span and its cokernel is zero.

Projecting every column of the full matrices gives
\[
 F_V=\begin{pmatrix}0&0&0&0\\
 [7]-[3]&0&0&0\\{}[8]&0&0&0\\0&[6]&-[3]&0
 \end{pmatrix},\qquad F_W=0,
\]
and the original Cartan ratios
$\operatorname{diag}(q^9,q^7,q^7,q^5)$ and $q^3I$.
The full divided square has the actual entry
\[
 (F_V^{(2)})_{144,126}
 =\frac{[6]([7]-[3])-[3][8]}{[2]}
 =q^{10}-q^4-q^{-4}+q^{-10}.
\]
The relation with the earlier restricted raw tableau lattice
is especially explicit. Its four raw words have images
$(e_0,e_1,-pe_2,-pe_3)$, and therefore their inclusion matrix
in $M_A$ is $\operatorname{diag}(1,1,-p,-p)$, with the
previously proved cokernel $(A/(p))^2$. Enlarging the raw
domain to $Y_A$ supplies other original words. The section
chosen above gives the following four preimages, each with
coefficient $+1$:
\[
\begin{array}{c|l}
T&(123)(123)(124)(12)(12)(1)(1)\\
T_1&(123)(123)(124)(12)(23)(1)(1)\\
T_2&(123)(123)(123)(12)(12)(2)(3)\\
Q&(123)(123)(124)(12)(13)(2)(3).
\end{array}
\]
For the last two rows the active $V$-string is $C_{12}=(2)(3)$;
the middle formula in \eqref{eq:full-contextual-relations}
gives $(2)(3)=-p^{-1}(4)(1)$ with these exact contexts.
This proves their stated images directly. The $T_1$ and
$Q$ column-order adjustments use the integral $(2,2)$
commutations already listed.

Put $Y_{\mathcal D,A}=\pi_A^{-1}
(\bigoplus_{n\in\mathcal D}Ab_n)$. Inclusion of the four
old raw words into this module and the composite
$Y_{\mathcal D,A}\xrightarrow{\pi_A}
\bigoplus_{n\in\mathcal D}Ab_n\twoheadrightarrow M_A$
give the commutative square
\[
\begin{array}{ccc}
A^4_{\rm old\ raw}&\longrightarrow&Y_{\mathcal D,A}\\
\big\downarrow{\scriptstyle\operatorname{diag}(1,1,-p,-p)}
 &&\big\downarrow{\scriptstyle\pi_A\ {\rm followed\ by\ quotient}}\\
M_A&\xrightarrow{\quad I\quad}&M_A.
\end{array}
\]
Its right arrow is onto by the four additional preimages.
Its left arrow retains its exact nonzero cokernel; enlarging
the domain does not erase that restricted-lattice defect.
Thus the earlier four-vector construction is recovered by
an exact morphism from the full original generator module.
Its integral endpoint map, conductor factorization and
specialization kernel/cokernel remain the maps already
proved in Section~\ref{sec:gct-generator-interval}.
The extra $1256$ basis vectors here are original source
vectors with their original raising and lowering actions;
they are not chosen to repair the trace of a hypothetical
four-dimensional quantum-group representation.
